    {{ waterMarkText }}

    \renewcommand{\DTRPcs}{ThreePhaseFault} % DTR pcs definition
    \renewcommand{\OCname}{TransientBoltedConsumption}


    \subsection{Test I5 - Transient fault ride-through in consumption}

    Checks for compliant behavior of the BESS under a grid scenario where there
    is a symmetric three-phase fault at one of the four transmission lines
    (at a 1\% distance from the PDR). The BESS should remain connected and
    be able to supply the necessary reactive injection.

    \GridCircuitZthree

    \noindent Important reminders:
    \begin{itemize}
        \item The time for fault clearance, $T_{clear}$, is a parameter that
        is specified in the DTR and depends on the voltage level in
        question:
        \begin{itemize}
            \item HTB3 (400 kV): $T_{clear}$ = 85 ms
            \item HTB2 (150 kV and 225 kV): $T_{clear}$ = 85 ms
            \item HTB1 (63 kV and 90 kV): $T_{clear}$ = 150 ms
        \end{itemize}
    \end{itemize}

    In \Dynawo this fault has been modeled by splitting Line4 into Line4a
    and Line4b, and then inserting an intermediate bus between them, to
    which a NodeFault model is attached.

    Note that although the DTR PCS also requires simulating a one-phase
    fault, this tool will not implement it, as it would require the
    simulation of a three-phase model. All simulations are only
    positive-sequence.

    \begin{description}
        \item Initial conditions used at the PDR bus:
        \begin{itemize}
            \item $P = P_{max\_consumption\_unite}$
            \item $Q = 0$
            \item $U = U_{dim}$
        \end{itemize}
    \end{description}

    \subsubsection{Simulation parameters}

    Solver and parameters used in the simulation:
    \begin{center}
        \begin{tabular}{lc}
            \toprule
           \textbf{Parameter} & \textbf{Value (default)} \\
            \midrule
            \BLOCK{for row in solverPCSI16z3ThreePhaseFaultTransientBoltedConsumption}
            {{row[0]}}         & {{row[1]}}                         \\
            \BLOCK{endfor}
            \bottomrule
        \end{tabular}
    \end{center}

    \GridSimulationHeaderZthree{
        \GridCurvesDescriptionZthree
        \GridCurvesItZthreeDescriptionZthree
        \GridCurvesTapZthreeDescriptionZthree
    }
    \GridCurvesZthree
    \GridCurvesItZthree
    \GridCurvesTapZthree
    \\[2\baselineskip]
    Go to  {{ linkPCSI16z3ThreePhaseFaultTransientBoltedConsumption }}

    \subsubsection{Analysis of results}

    \noindent Analysis of results for PCS \DTRPcs. Values extracted
    from the simulation:

    \begin{center}
        \scriptsize
        \begin{tabular}{@{}lccccccccc@{}}
            \toprule
            & \multicolumn{3}{c}{Pre-event} & \multicolumn{3}{c}{Event} & \multicolumn{3}{c}{Pos-Event} \\
            \cmidrule(lr){2-4}\cmidrule(lr){5-7}\cmidrule(lr){8-10}
            & {MXE}      & {ME}       & {MAE}      & {MXE}      & {ME}       & {MAE}      & {MXE}      & {ME}       & {MAE}      \\
            \midrule
            \BLOCK{for row in rmPCSI16z3ThreePhaseFaultTransientBoltedConsumption}
            {{row[0]}} & {{row[1]}} & {{row[2]}} & {{row[3]}} & {{row[4]}} & {{row[5]}} & {{row[6]}} & {{row[7]}} & {{row[8]}} & {{row[9]}} \\
            \BLOCK{endfor}
            \bottomrule
        \end{tabular}
    \end{center}

    \subsubsection{Compliance checks}

    {{ missedColumns }}

    \noindent Compliance thresholds on the curves:
    \begin{center}
        \scriptsize
        \begin{tabular}{@{}lccccccccc@{}}
            \toprule
            & \multicolumn{3}{c}{Pre-event} & \multicolumn{3}{c}{Event} & \multicolumn{3}{c}{Pos-Event} \\
            \cmidrule(lr){2-4}\cmidrule(lr){5-7}\cmidrule(lr){8-10}
            & {MXE}      & {ME}       & {MAE}      & {MXE}      & {ME}       & {MAE}      & {MXE}      & {ME}       & {MAE}      \\
            \midrule
            \BLOCK{for row in thmPCSI16z3ThreePhaseFaultTransientBoltedConsumption}
            {{row[0]}} & {{row[1]}} & {{row[2]}} & {{row[3]}} & {{row[4]}} & {{row[5]}} & {{row[6]}} & {{row[7]}} & {{row[8]}} & {{row[9]}} \\
            \BLOCK{endfor}
            \bottomrule
        \end{tabular}
    \end{center}

    \noindent Compliance checks on the curves:
    \begin{center}
        \scriptsize
        \begin{tabular}{@{}lcccccccccc@{}}
            \toprule
            & \multicolumn{3}{c}{Pre-event} & \multicolumn{3}{c}{Event} & \multicolumn{3}{c}{Pos-Event} & \\
            \cmidrule(lr){2-4}\cmidrule(lr){5-7}\cmidrule(lr){8-10}
            & {MXE}      & {ME}       & {MAE}      & {MXE}      & {ME}       & {MAE}      & {MXE}      & {ME}       & {MAE}      & Compl.      \\
            \midrule
            \BLOCK{for row in emPCSI16z3ThreePhaseFaultTransientBoltedConsumption}
            {{row[0]}} & {{row[1]}} & {{row[2]}} & {{row[3]}} & {{row[4]}} & {{row[5]}} & {{row[6]}} & {{row[7]}} & {{row[8]}} & {{row[9]}} & {{row[10]}} \\
            \BLOCK{endfor}
            \bottomrule
        \end{tabular}
    \end{center}

    \noindent Compliance checks on the Active Power Recovery:
    \begin{center}
        \scriptsize
        \begin{tabular}{cllc}
            \toprule
            Variable   & Error      & Threshold   & Check      \\
            \midrule
            \BLOCK{for row in aprPCSI16z3ThreePhaseFaultTransientBoltedConsumption}
            {{row[0]}} & {{row[1]}} & {{row[2]}}  & {{row[3]}} \\
            \BLOCK{endfor}
            \bottomrule
        \end{tabular}
    \end{center}

    \noindent In steady state after the event, the absolute average error must not exceed {{steadystatethreshold}}\% (configured value):
    \begin{center}
        \scriptsize
        \begin{tabular}{cllc}
            \toprule
            Variable   & MAE        & Final Error & Compliance \\
            \midrule
            \BLOCK{for row in ssemPCSI16z3ThreePhaseFaultTransientBoltedConsumption}
            {{row[0]}} & {{row[1]}} & {{row[2]}}  & {{row[3]}} \\
            \BLOCK{endfor}
            \bottomrule
        \end{tabular}
    \end{center}
